\medskip
\hypertarget{native-109-article080}{}\pdfbookmark[2]{80.}{article-080}
\begin{center}80.\end{center}
Ist der Winkel $\Pi$ gegeben, so bestimmen sich $p$ und $e$ aus den Gleichungen
\[
\begin{aligned}
p&=\frac{rr'(\cos(N-\Pi)-\cos(N'-\Pi))}{r\cos(N-\Pi)-r'\cos(N'-\Pi)}\\[4pt]
e&=\frac{r'-r}{r\cos(N-\Pi)-r'\cos(N'-\Pi)}\mathpunct{\cdot}
\end{aligned}
\]
Der gemeinsame Nenner dieser Formeln lässt sich in die Form $a\cos(A-\Pi)$ bringen, wobei $a$ und $A$ von $\Pi$ unabhängig sind. Bezeichnet nämlich $H$ einen beliebigen Winkel, so gilt
\[
\begin{aligned}
&r\cos(N-\Pi)-r'\cos(N'-\Pi)\\
&\quad=(r\cos(N-H)-r'\cos(N'-H))\cos(H-\Pi)\\
&\qquad-(r\sin(N-H)-r'\sin(N'-H))\sin(H-\Pi),
\end{aligned}
\]
und der Nenner ist somit $=a\cos(A-\Pi)$, wenn $a$ und $A$ aus den Gleichungen bestimmt werden
\SourceLeaf{110}
\[
\begin{aligned}
r\cos(N-H)-r'\cos(N'-H)&=a\cos(A-H)\\
r\sin(N-H)-r'\sin(N'-H)&=a\sin(A-H).
\end{aligned}
\]
So ergibt sich
\[
\begin{aligned}
p&=\frac{2rr'\sin\tfrac12(N'-N)\sin(\tfrac12N+\tfrac12N'-\Pi)}{a\cos(A-\Pi)}\\[4pt]
e&=\frac{r'-r}{a\cos(A-\Pi)}\mathpunct{\cdot}
\end{aligned}
\]
Diese Formeln sind besonders zweckmäßig, wenn $p$ und $e$ für mehrere Werte von $\Pi$ bei unveränderten $r$, $r'$, $N$, $N'$ berechnet werden sollen. — Bei der Berechnung der Hilfsgrößen $a$ und $A$ kann der Winkel $H$ frei gewählt werden. Daher setzt man zweckmäßigerweise $H=\tfrac12(N+N')$. Die Gleichungen für diese Größen gehen dann über in
\[
\begin{aligned}
(r'-r)\cos\tfrac12(N'-N)&=-a\cos(A-\tfrac12N-\tfrac12N')\\
(r'+r)\sin\tfrac12(N'-N)&=-a\sin(A-\tfrac12N-\tfrac12N').
\end{aligned}
\]
Hat man also den Winkel $A$ aus der Gleichung
\[
\operatorname{tang}(A-\tfrac12N-\tfrac12N')=\frac{r'+r}{r'-r}\operatorname{tang}\tfrac12(N'-N),
\]
bestimmt, so erhält man unmittelbar
\[
e=-\frac{\cos(A-\tfrac12N-\tfrac12N')}{\cos\tfrac12(N'-N)\cos(A-\Pi)}\mathpunct{\cdot}
\]
Die Berechnung des Logarithmus von $\dfrac{r'+r}{r'-r}$ lässt sich durch das schon mehrfach erläuterte Verfahren abkürzen.
